\documentclass[aspectratio=169]{beamer}

\usetheme{Madrid}
\usepackage[spanish]{babel}
\usepackage[utf8]{inputenc}
\usepackage{textcomp}
\usepackage[T1]{fontenc}
\usepackage{lmodern}
\usepackage{graphicx}
\usepackage{booktabs}
\usepackage{verbatim}
\usepackage{amssymb}
\usepackage{pifont}
\usepackage{moreverb}
\usepackage[most]{tcolorbox}
\usepackage{tikz}
\usetikzlibrary{shapes, arrows, positioning}

% Configuración de TikZ
\tikzstyle{block} = [
  rectangle, draw, rounded corners,
  text centered, minimum height=1cm,
  minimum width=2.5cm, font=\small, fill=gray!5
]
\tikzstyle{arrow} = [->, thick]

% Configuración de tcolorbox
\tcbset{
  colback=gray!5, colframe=gray!50,
  boxrule=0.4pt, sharp corners,
}

\title[BookMate]{BookMate — Sistema de Recomendación de Libros Académicos}
\subtitle{Diseño Arquitectónico: Stack Tecnológico, Arquitectura C4 y Patrones de Diseño}
\author[GRUPO 6]{Delgado R., G. \and Osorio M., A. \and Rojas A., J. \and Torres R., J. \and Valverde G., Y. \and Villanueva A., F.}
\institute[FC-UNI]{UNIVERSIDAD NACIONAL DE INGENIERÍA}
\titlegraphic{\includegraphics[height=1.4cm]{example-image}}
% NOTA: Reemplazar 'example-image' con 'Imagenes/logo_uni.png' cuando esté disponible

\date{Noviembre 2025}

\begin{document}

% Portada
\begin{frame}
\titlepage
\end{frame}

% Agenda
\begin{frame}
\frametitle{Agenda}
\tableofcontents
\end{frame}

\section{Introducción al Diseño Arquitectónico}

\begin{frame}
\frametitle{Fase de Diseño Arquitectónico}

\begin{columns}
\begin{column}{0.5\textwidth}
\textbf{Objetivo de esta Fase:}
\begin{itemize}
\item Definir CÓMO se implementará el sistema
\item Seleccionar tecnologías específicas
\item Diseñar arquitectura técnica detallada
\item Aplicar patrones de diseño
\item Documentar decisiones arquitectónicas
\end{itemize}

\vspace{0.3cm}

\textbf{Enfoque:}
\begin{tcolorbox}[colback=green!5, colframe=green!50]
\textbf{Con tecnologías específicas}\\
Spring Boot, PostgreSQL, Python, Docker, etc.
\end{tcolorbox}
\end{column}

\begin{column}{0.5\textwidth}
\textbf{Entregables de Semana 13:}
\begin{enumerate}
\item Prototipo navegable actualizado
\item Plan de desarrollo del software
\item Documento de arquitectura completo
\item Arquitectura C4 Niveles 3-4
\item Presentación Beamer
\end{enumerate}
\end{column}
\end{columns}

\vspace{0.3cm}

\begin{tcolorbox}[colback=yellow!5, colframe=yellow!50]
\textbf{Nota:} A diferencia de la fase de análisis (Semana 11), aquí definimos tecnologías concretas, clases de implementación y arquitectura física.
\end{tcolorbox}

\end{frame}

\begin{frame}
\frametitle{Transición: Análisis → Diseño}

\begin{columns}
\begin{column}{0.5\textwidth}
\textbf{Fase de Análisis (S11):}
\begin{itemize}
\item Modelo conceptual (BCE)
\item Entidades del dominio
\item Arquitectura conceptual
\item Sin tecnologías específicas
\end{itemize}
\end{column}

\begin{column}{0.5\textwidth}
\textbf{Fase de Diseño (S13):}
\begin{itemize}
\item Clases de implementación (JPA)
\item Stack tecnológico definido
\item Arquitectura física
\item Patrones de diseño aplicados
\end{itemize}
\end{column}
\end{columns}

\vspace{0.5cm}

\begin{center}
\begin{tikzpicture}[node distance=3cm]
\node[block, fill=LightBlue] (analisis) {Análisis\\Conceptual};
\node[block, fill=LightGreen, right=of analisis] (diseno) {Diseño\\Arquitectónico};
\draw[arrow] (analisis) -- node[above] {Transformación} (diseno);
\end{tikzpicture}
\end{center}

\end{frame}

\section{Stack Tecnológico}

\begin{frame}
\frametitle{Stack Tecnológico Completo}

\begin{table}
\centering
\tiny
\begin{tabular}{p{2.5cm}p{2.5cm}p{3cm}p{3.5cm}}
\toprule
\textbf{Componente} & \textbf{Tecnología} & \textbf{Versión} & \textbf{Propósito} \\
\midrule
\textbf{Backend} & Spring Boot & 3.2.0 & Framework principal \\
& Spring Data JPA & (incluido) & Acceso a datos \\
& Spring Security & (incluido) & Autenticación \\
& Flyway & 9.x & Migraciones BD \\
\midrule
\textbf{Base de Datos} & PostgreSQL & 16 & Persistencia \\
& HikariCP & (incluido) & Connection pool \\
\midrule
\textbf{Servicio IA} & Python & 3.11+ & Lenguaje IA \\
& Flask & 2.3+ & Framework web \\
& Sentence Transformers & 2.2+ & Embeddings \\
& NumPy/SciPy & 1.24+/1.11+ & Cálculos numéricos \\
\midrule
\textbf{Frontend} & HTML5/CSS3/JS & ES6+ & Interfaz \\
& Bootstrap & 5.1.3 & UI framework \\
\midrule
\textbf{Containerización} & Docker & 24+ & Contenedores \\
& Docker Compose & 2.20+ & Orquestación \\
\bottomrule
\end{tabular}
\end{table}

\end{frame}

\begin{frame}
\frametitle{Justificación de Tecnologías Principales}

\textbf{Spring Boot:}
\begin{itemize}
\item Convención sobre configuración (reduce código boilerplate)
\item Ecosistema completo (JPA, Security, Web)
\item Excelente documentación y comunidad
\item Ideal para APIs REST
\end{itemize}

\textbf{PostgreSQL:}
\begin{itemize}
\item Open source y robusto
\item Excelente soporte ACID
\item Compatible con Spring Data JPA
\item Extensible y ampliamente usado
\end{itemize}

\textbf{Python para IA:}
\begin{itemize}
\item Ecosistema rico en ML/NLP
\item Sentence Transformers fácil de usar
\item Modelos preentrenados disponibles
\item Ideal para microservicios de IA
\end{itemize}

\end{frame}

\section{Arquitectura C4 Completa}

\begin{frame}
\frametitle{C4 Model - Nivel 1: Contexto}

\begin{center}
\includegraphics[width=0.9\textwidth]{example-image}
% NOTA: Reemplazar con diagrama C4 Nivel 1 (Contexto)
% Mostrar:
% - Personas: Usuario, Administrador
% - Sistema BookMate (centro)
% - Sistema de IA (externo)
% - Relaciones entre ellos
% - Tecnologías: Spring Boot, PostgreSQL, Python
\end{center}

\textbf{Elementos:}
\begin{itemize}
\item \textbf{Usuarios:} Estudiante, profesor, investigador
\item \textbf{Administradores:} Bibliotecario, gestor de catálogo
\item \textbf{Sistema BookMate:} Backend Spring Boot + Frontend
\item \textbf{Servicio de IA:} Microservicio Python con Sentence Transformers
\end{itemize}

\end{frame}

\begin{frame}
\frametitle{C4 Model - Nivel 2: Contenedores}

\begin{center}
\includegraphics[width=0.85\textwidth]{example-image-a}
% NOTA: Reemplazar con diagrama C4 Nivel 2 (Contenedores)
% Mostrar:
% - Aplicación Web (HTML/CSS/JS)
% - Sistema de Gestión (Spring Boot)
% - Sistema de Almacenamiento (PostgreSQL)
% - Servicio de IA (Python Flask)
% - Flujos de comunicación con tecnologías específicas
\end{center}

\textbf{Contenedores:}
\begin{itemize}
\item \textbf{Aplicación Web:} HTML5, CSS3, JavaScript, Bootstrap
\item \textbf{Sistema de Gestión:} Spring Boot 3.2.0, Java 17+
\item \textbf{Sistema de Almacenamiento:} PostgreSQL 16
\item \textbf{Servicio de IA:} Python 3.11, Flask 2.3, Sentence Transformers
\end{itemize}

\end{frame}

\begin{frame}
\frametitle{C4 Model - Nivel 3: Componentes}

\begin{center}
\includegraphics[width=0.8\textwidth]{example-image-b}
% NOTA: Reemplazar con diagrama C4 Nivel 3 (Componentes)
% Mostrar componentes del backend:
% - Controllers (LibrosController, AutorController, RecommendationController)
% - Services (LibrosService, RecommendationService, AIServiceClient)
% - Repositories (LibrosRepository, AutorRepository, EmbeddingRepository)
% - DTOs y Entities
% - Relaciones entre componentes
\end{center}

\textbf{Componentes Principales:}
\begin{itemize}
\item \textbf{Controllers:} LibrosController, AutorController, RecommendationController
\item \textbf{Services:} LibrosService, RecommendationService, AIServiceClient, HeuristicRecommender
\item \textbf{Repositories:} LibrosRepository, AutorRepository, EmbeddingRepository
\item \textbf{DTOs:} LibroDTO, AutorDTO, RecomendacionDTO
\end{itemize}

\end{frame}

\begin{frame}
\frametitle{C4 Model - Nivel 4: Código}

\begin{center}
\includegraphics[width=0.75\textwidth]{example-image-c}
% NOTA: Reemplazar con diagrama C4 Nivel 4 (Código)
% Mostrar métodos específicos:
% - RecommendationController.obtenerRecomendaciones(bookId)
% - RecommendationService.obtenerRecomendaciones(bookId)
% - AIServiceClient.calcularSimilitud(bookId)
% - HeuristicRecommender.recomendar(bookId)
% - Flujo de llamadas entre métodos
\end{center}

\textbf{Métodos Clave:}
\begin{itemize}
\item \textbf{RecommendationController:} obtenerRecomendaciones(bookId)
\item \textbf{RecommendationService:} obtenerRecomendaciones(bookId), decidirMetodo()
\item \textbf{AIServiceClient:} calcularSimilitud(bookId), verificarDisponibilidad()
\item \textbf{HeuristicRecommender:} recomendar(bookId), calcularPuntuacion()
\end{itemize}

\end{frame}

\section{Patrones de Diseño}

\begin{frame}
\frametitle{Patrón Strategy: Recomendaciones}

\textbf{Aplicación:} Sistema de recomendaciones (IA vs Heurísticas)

\textbf{Implementación:}
\begin{itemize}
\item \textbf{Interfaz:} RecommenderStrategy
\item \textbf{Implementaciones:} AIRecommender, HeuristicRecommender
\item \textbf{Contexto:} RecommendationService decide qué estrategia usar
\end{itemize}

\textbf{Ventajas:}
\begin{itemize}
\item Facilita agregar nuevos algoritmos
\item Permite cambiar estrategia sin modificar código cliente
\item Fallback automático a heurísticas si IA falla
\end{itemize}

\vspace{0.3cm}

\begin{center}
\includegraphics[width=0.6\textwidth]{example-image}
% NOTA: Reemplazar con diagrama del patrón Strategy
% Mostrar: RecommendationService → AIRecommender / HeuristicRecommender
\end{center}

\end{frame}

\begin{frame}
\frametitle{Patrón Repository: Acceso a Datos}

\textbf{Aplicación:} Abstracción de acceso a base de datos

\textbf{Implementación:}
\begin{itemize}
\item Interfaces Spring Data JPA
\item Implementación automática por Spring
\item Consultas derivadas de métodos
\end{itemize}

\textbf{Ejemplo:}
\begin{verbatim}
public interface LibrosRepository 
    extends JpaRepository<Libro, Long> {
    List<Libro> findByTituloContaining(String titulo);
    List<Libro> findByAutorId(Long autorId);
    Page<Libro> findAll(Pageable pageable);
}
\end{verbatim}

\textbf{Ventajas:}
\begin{itemize}
\item Separa lógica de negocio de persistencia
\item Facilita testing (mocks)
\item Reduce código boilerplate
\end{itemize}

\end{frame}

\begin{frame}
\frametitle{Patrón Service Layer: Lógica de Negocio}

\textbf{Aplicación:} Capa intermedia entre controllers y repositories

\textbf{Responsabilidades:}
\begin{itemize}
\item Lógica de negocio
\item Coordinación entre componentes
\item Transformación Entidad → DTO
\item Aplicación de reglas de negocio
\end{itemize}

\textbf{Ejemplo:}
\begin{verbatim}
@Service
public class RecommendationService {
    public List<LibroDTO> obtenerRecomendaciones(Long bookId) {
        // Lógica de negocio
        // Decisión de método
        // Transformación a DTOs
    }
}
\end{verbatim}

\textbf{Ventajas:}
\begin{itemize}
\item Separa lógica de negocio de presentación
\item Facilita reutilización
\item Facilita testing
\end{itemize}

\end{frame}

\begin{frame}
\frametitle{Patrón DTO: Transferencia de Datos}

\textbf{Aplicación:} Transferencia de datos entre capas

\textbf{Propósito:}
\begin{itemize}
\item Controlar qué datos se exponen
\item Evitar problemas de lazy loading
\item Facilitar versionado de API
\end{itemize}

\textbf{Ejemplo:}
\begin{verbatim}
public class LibroDTO {
    private Long id;
    private String titulo;
    private String autorNombre;
    // Solo campos necesarios
}
\end{verbatim}

\textbf{Ventajas:}
\begin{itemize}
\item Control de exposición de datos
\item Evita problemas de serialización
\item Facilita evolución de API
\end{itemize}

\end{frame}

\section{Decisiones Arquitectónicas}

\begin{frame}
\frametitle{Decisión 1: Microservicio Python Separado}

\textbf{Decisión:} Implementar servicio de IA como microservicio Python separado

\textbf{Justificación:}
\begin{itemize}
\item Separación de responsabilidades (IA es dominio diferente)
\item Escalabilidad independiente
\item Tecnología adecuada (Python mejor para ML/NLP)
\item Aislamiento de fallos (fallback a heurísticas)
\end{itemize}

\textbf{Trade-offs:}
\begin{itemize}
\item ✅ Ventajas: Separación clara, escalabilidad, tecnología adecuada
\item ❌ Desventajas: Complejidad de despliegue, latencia de red
\end{itemize}

\vspace{0.3cm}

\begin{center}
\includegraphics[width=0.6\textwidth]{example-image-a}
% NOTA: Reemplazar con diagrama de microservicios
% Mostrar: Backend Spring Boot ↔ Servicio Python IA
\end{center}

\end{frame}

\begin{frame}
\frametitle{Decisión 2: Fallback Automático a Heurísticas}

\textbf{Decisión:} Si servicio IA no disponible, usar automáticamente heurísticas

\textbf{Justificación:}
\begin{itemize}
\item Alta disponibilidad (sistema siempre retorna recomendaciones)
\item Resiliencia (no depende completamente de servicio externo)
\item Experiencia de usuario (siempre recibe resultados)
\end{itemize}

\textbf{Implementación:}
\begin{itemize}
\item Strategy pattern para seleccionar método
\item Health check del servicio IA
\item Fallback transparente
\end{itemize}

\vspace{0.3cm}

\begin{tcolorbox}[colback=green!5, colframe=green!50]
\textbf{Resultado:} Sistema siempre funcional, incluso si servicio IA falla.
\end{tcolorbox}

\end{frame}

\begin{frame}
\frametitle{Decisión 3: Flyway para Migraciones}

\textbf{Decisión:} Usar Flyway para versionado y migración de esquema de BD

\textbf{Justificación:}
\begin{itemize}
\item Control de versiones del esquema
\item Migraciones automáticas al iniciar aplicación
\item Historial completo de cambios
\item Rollback de migraciones
\end{itemize}

\textbf{Estructura:}
\begin{verbatim}
src/main/resources/db/migration/
├── V1__Initial_schema.sql
├── V2__Add_authors_table.sql
└── V3__Add_embeddings_table.sql
\end{verbatim}

\textbf{Ventajas:}
\begin{itemize}
\item Reproducibilidad (mismo esquema en todos los entornos)
\item Control de cambios
\item Integración nativa con Spring Boot
\end{itemize}

\end{frame}

\begin{frame}
\frametitle{Decisión 4: Docker Compose para Orquestación}

\textbf{Decisión:} Usar Docker Compose para orquestar los 3 servicios

\textbf{Justificación:}
\begin{itemize}
\item Desarrollo local fácil (un comando para levantar todo)
\item Reproducibilidad (mismo entorno en todos los desarrolladores)
\item Aislamiento (cada servicio en su contenedor)
\item Configuración declarativa (docker-compose.yml documenta arquitectura)
\end{itemize}

\textbf{Servicios:}
\begin{itemize}
\item Backend Spring Boot (puerto 8080)
\item PostgreSQL (puerto 5432)
\item Servicio Python IA (puerto 5000)
\end{itemize}

\vspace{0.3cm}

\begin{center}
\includegraphics[width=0.7\textwidth]{example-image-b}
% NOTA: Reemplazar con diagrama Docker Compose
% Mostrar: 3 contenedores, red, volúmenes
\end{center}

\end{frame}

\section{Comunicación Entre Servicios}

\begin{frame}
\frametitle{Backend ↔ Frontend}

\textbf{Protocolo:} HTTP REST  
\textbf{Formato:} JSON  
\textbf{Autenticación:} JWT tokens

\textbf{Endpoints principales:}
\begin{itemize}
\item \texttt{GET /api/books} - Listar libros
\item \texttt{GET /api/books/\{id\}} - Detalle de libro
\item \texttt{GET /api/recommendations/\{bookId\}} - Recomendaciones
\item \texttt{POST /api/auth/login} - Autenticación
\end{itemize}

\textbf{Ejemplo de petición:}
\begin{verbatim}
GET /api/recommendations/123 HTTP/1.1
Host: localhost:8080
Authorization: Bearer {jwt_token}
\end{verbatim}

\end{frame}

\begin{frame}
\frametitle{Backend ↔ Servicio IA}

\textbf{Protocolo:} HTTP REST  
\textbf{Formato:} JSON  
\textbf{Timeout:} 5 segundos  
\textbf{Retry:} 1 intento

\textbf{Endpoints del servicio IA:}
\begin{itemize}
\item \texttt{POST /api/embeddings/generate} - Generar embedding
\item \texttt{GET /api/similarity/\{bookId\}} - Calcular similitud
\item \texttt{GET /health} - Health check
\end{itemize}

\textbf{Manejo de errores:}
\begin{itemize}
\item Timeout: Usar heurísticas
\item Error de conexión: Usar heurísticas
\item Fallback automático transparente
\end{itemize}

\vspace{0.3cm}

\begin{center}
\includegraphics[width=0.7\textwidth]{example-image-c}
% NOTA: Reemplazar con diagrama de flujo de comunicación
% Mostrar: Backend → IA → Backend, con manejo de errores
\end{center}

\end{frame}

\section{Seguridad}

\begin{frame}
\frametitle{Autenticación y Autorización}

\textbf{Tecnología:} Spring Security + JWT

\textbf{Flujo:}
\begin{enumerate}
\item Usuario envía credenciales (\texttt{POST /api/auth/login})
\item Backend valida credenciales
\item Backend genera JWT token
\item Frontend almacena token
\item Frontend envía token en headers
\item Backend valida token en cada petición
\end{enumerate}

\textbf{Roles:}
\begin{itemize}
\item \textbf{USER:} Usuario estándar (buscar, ver detalles, recomendaciones)
\item \textbf{ADMIN:} Administrador (CRUD de libros y autores)
\end{itemize}

\textbf{Implementación:}
\begin{itemize}
\item \texttt{@PreAuthorize("hasRole('ADMIN')")} en métodos
\item Filtros de Spring Security
\item Tokens JWT con expiración
\end{itemize}

\end{frame}

\begin{frame}
\frametitle{Validación de Entrada}

\textbf{Tecnología:} Bean Validation (\texttt{@Valid}, \texttt{@NotNull}, \texttt{@Size})

\textbf{Ejemplo:}
\begin{verbatim}
public class CrearLibroDTO {
    @NotBlank(message = "Título es obligatorio")
    @Size(max = 200)
    private String titulo;
    
    @NotNull
    @DecimalMin(value = "0.0")
    private BigDecimal precio;
}
\end{verbatim}

\textbf{Protección:}
\begin{itemize}
\item SQL Injection: Prevenido por JPA (prepared statements)
\item XSS: Validación de entrada + escapado HTML
\item CSRF: Tokens JWT + CORS configurado
\end{itemize}

\end{frame}

\section{Resumen y Conclusiones}

\begin{frame}
\frametitle{Resumen del Diseño Arquitectónico}

\begin{columns}
\begin{column}{0.5\textwidth}
\textbf{Stack Tecnológico:}
\begin{itemize}
\item Backend: Spring Boot 3.2.0
\item Base de datos: PostgreSQL 16
\item Servicio IA: Python 3.11 + Flask
\item Frontend: HTML5/CSS3/JS
\item Containerización: Docker Compose
\end{itemize}
\end{column}

\begin{column}{0.5\textwidth}
\textbf{Arquitectura:}
\begin{itemize}
\item Microservicios ligero (2 servicios)
\item Arquitectura C4 completa (4 niveles)
\item Patrones: Strategy, Repository, Service Layer, DTO
\item Fallback automático a heurísticas
\item Alta disponibilidad y resiliencia
\end{itemize}
\end{column}
\end{columns}

\vspace{0.5cm}

\begin{tcolorbox}[colback=blue!5, colframe=blue!50, title=Objetivo Cumplido]
Hemos definido CÓMO se implementará el sistema con tecnologías específicas, arquitectura detallada y patrones de diseño aplicados. Esto proporciona una base sólida para el diseño detallado (Semana 15).
\end{tcolorbox}

\end{frame}

\begin{frame}
\frametitle{Logros del Diseño Arquitectónico}

\begin{enumerate}
\item ✅ \textbf{Stack tecnológico} definido y justificado
\item ✅ \textbf{Arquitectura C4} niveles 1-4 documentada
\item ✅ \textbf{Patrones de diseño} identificados y aplicados
\item ✅ \textbf{Decisiones arquitectónicas} documentadas
\item ✅ \textbf{Comunicación entre servicios} especificada
\item ✅ \textbf{Seguridad} planificada (JWT, validación)
\item ✅ \textbf{Containerización} con Docker Compose
\item ✅ \textbf{Plan de desarrollo} actualizado
\end{enumerate}

\vspace{0.5cm}

\begin{tcolorbox}[colback=green!5, colframe=green!50, title=Próximos Pasos]
Semana 15: Diseño detallado con clases de implementación, diagramas de secuencia, componentes y despliegue.
\end{tcolorbox}

\end{frame}

\begin{frame}
\frametitle{¡Preguntas?}

\begin{center}
\Large ¡Gracias por su atención! \\

\vspace{1cm}

\textbf{¿Preguntas sobre el diseño arquitectónico?}

\vspace{1cm}

\normalsize
\textbf{Grupo 6}\\
Universidad Nacional de Ingeniería\\
CC341 - Ingeniería de Software\\
Ciclo 2025-2

\vspace{0.5cm}

\textbf{Próxima entrega:} Semana 15 - Diseño Detallado (clases de implementación)

\end{center}

\end{frame}

\end{document}

